% Responsable: Persona 3. ~100 palabras + tabla.
\section{Herramienta para el cliente}

La tabla~\ref{tab:comparativa} resume las cuatro opciones. Las columnas de
distancias, k-means y vecinos se calcularon con las cabeceras municipales
(Paso~5); k-means corresponde a $k = 20$ (mediana de diez semillas).

\begin{table}[ht]
  \centering
  \footnotesize
  \setlength{\tabcolsep}{4pt}
  \begin{tabular}{lrrrrrrr}
\toprule
 & & Error máx. & \multicolumn{3}{c}{Error en distancias (\%)} & k-means & 10 vecinos \\
\cmidrule(lr){4-6}
Proyección & $\delta$ & escala (\%) & mediana & P95 & máx. & cambios (\%) & cambian (\%) \\
\midrule
Mercator & 0.1402 & 7.26 & 6.07 & 10.21 & 18.04 & 0.67 & 0.69 \\
Lambert INEGI & 0.0132 & 0.66 & 0.28 & 0.49 & 0.60 & 0.24 & 0.08 \\
Óptima (México) & 0.0018 & 0.09 & 0.04 & 0.13 & 1.07 & 0.02 & 0.12 \\
Óptima (envolvente) & 0.0039 & 0.19 & 0.07 & 0.16 & 0.19 & 0.02 & 0.08 \\
\bottomrule
\end{tabular}

  \caption{Comparación de las proyecciones sobre México. El error máximo de
    escala supone el mapa reescalado para centrar $\sigma$ en 1. El máximo de
    la óptima (México) viene de pares que cruzan el Golfo de California.
    Fuente: elaboración propia con datos del INEGI.}
  \label{tab:comparativa}
\end{table}

Para convertir coordenadas, el repositorio incluye la función
\texttt{latlon\_a\_xy(lat, lon, proyeccion="recomendada")} en
\texttt{src/proyecciones/conversion.py}. Recibe grados (números, arreglos o
las columnas de un CSV) y devuelve $x, y$ en metros. La opción recomendada es
la óptima sobre la envolvente convexa, cuyo modelo ya ajustado viene en el
repositorio, así que no hace falta resolver nada; también acepta
\texttt{"optima"}, \texttt{"lambert"} y \texttt{"mercator"}; avisa si un
punto queda fuera de la región, y las filas sin coordenadas quedan vacías. Desde la terminal:
\begin{lstlisting}[basicstyle=\footnotesize\ttfamily]
python -m proyecciones.conversion clientes.csv clientes_xy.csv
\end{lstlisting}
Por ejemplo, con la proyección recomendada la distancia en el plano entre la
Ciudad de México y Monterrey es de 704.78~km, contra 705.93~km sobre la
esfera ($-0.16\,\%$), y entre Monterrey y Mérida, de 1212.99~km contra
1213.52~km ($-0.04\,\%$).
